%% 第四章 效用引导的渐进式遗忘

\chapter{效用引导的渐进式遗忘}
\chaptermark{效用引导的渐进式遗忘}

\section{引言}

第三章的事件关联图和结构化检索解决了"如何在长历史中找得更准"的问题。然而，长期部署中的记忆系统还面临另一个同样关键的挑战：\emph{持续写入新记忆必然导致存储的增长}。服务机器人每天产生海量的多模态感知数据：RGB 图像帧（单帧约 200 KB，以 5 fps 采集一小时即产生约 3.6 GB）、深度图、音频流、关节角度序列。在层级记忆树的 append-only 架构中，所有原始数据被保存在 L0 层，每个 L2 事件节点关联数帧至数十帧 L0 数据。同时，L1 场景图每帧包含 20+ 个物体实例和 30+ 条空间关系，L2–L4+ 层的 LLM 生成摘要随 episode 数量累积。不加选择的存储策略在长期部署中不可持续：存储成本随运行时间线性膨胀，同时每次 QA 都需要将完整的记忆树加载到内存中进行检索，过大的树结构也会拖慢图构建和语义搜索的速度。本文使用的评测缓存（仅存储树结构元数据和文本摘要，不含原始图像/音频文件）在 $|h|=100$ 时即已包含 60,836 个事件节点、68,214 个场景图和 577,052 条关系——这些小规模的树结构数据已经为"持续膨胀"提供了清晰的信号，而实际部署中叠加原始多模态文件后的存储压力将比这一数字高出数个数量级。

更关键的是，\emph{并非所有记忆都具有同等的保留价值}。这一观察不仅是工程直觉，在人类记忆研究中也有深刻的对应。认知心理学将人类长时记忆的巩固过程描述为一个高度选择性的机制：不是所有的经历都以同等的细节水平被保留。大脑在睡眠期间通过\emph{系统巩固}（systems consolidation）过程，将海马体中的情景记忆逐步重组并转移至新皮层，在这一过程中，与已有知识结构高度吻合的信息和带有强烈情感标记的事件被优先强化，而孤立的、重复性的日常细节则逐渐退化为要旨记忆（gist memory）——仅保留大致印象而丢失具体细节。这种"选择性保留+渐进式降级"的策略使人类能够在数十年的生命跨度中维持高效的记忆检索，而不被海量的日常细节淹没。

机器人的长期部署场景有着类似的需求。三天前的一次失败任务远比一周前的一次例行导航帧更值得保留；包含用户语音交互的事件（"make coffee, please""the bread is in the lower cabinet"）承载了任务的关键语义，类似于人类记忆中的情感标记事件；在事件图中作为结构枢纽的高连接度事件在多条检索路径中都起到桥梁作用，类似于被多个已有知识结构锚定的记忆。反之，大量高度重复的日常动作——反复的导航帧、连续的无对象静默帧——在信息内容上高度冗余，对问答的边际贡献极低，其工程角色类似于人类记忆中那些迅速退化为要旨的日常细节。

本章解决第一章所提出的第二个研究问题（RQ2）：如何在保持问答能力基本稳定的前提下，控制长期运行中记忆存储的持续膨胀。核心方法是\emph{效用引导的渐进式遗忘}（Utility-Guided Progressive Forgetting）：不是全或无地删除记忆，而是在三个层级上施加程度不等的信息削减——完整保留（Level 0）、去细节化（Level 1）、摘要保留（Level 2）——并根据综合新近性、语义独特性、任务重要性和图结构中心性四个维度的效用评分决定每个节点的遗忘等级。遗忘模块的目标不是提高 QA 准确率（不应期望"忘了冗余信息反而答得更准"），而是在保持问答稳定性的前提下控制存储开销。

本章 4.2 节介绍效用函数的三个分项及其计算方式。4.3 节介绍渐进式遗忘策略——三级操作的定义、遗忘豁免与安全机制、以及遗忘→图构建的先筛后连管线。4.4 节在 TEACh $|h|=50$ 上报告两档实验（Forget+Graph 保守压缩和 Level 2 强压缩）的存储压缩与答案一致性结果。4.5 节在 ARMARX 真实机器人数据上通过结构审计和局部细节切片补充图中心性引导遗忘的机制验证。4.6 节进行综合讨论。

\section{效用函数设计}

效用函数 $U(n)$ 对记忆树中每个 L2 事件节点 $n$ 计算一个 $[0, 1]$ 的连续值，表示该节点在后续问答中的预期价值。效用值越高，节点越应被保留；效用值越低，节点越可以被压缩。效用函数由三个分项加权组合而成：

\begin{equation}
U(n) = \alpha \cdot \text{Recency}(n) + \beta \cdot \text{Uniqueness}(n) + \gamma \cdot \text{Importance}(n)
\label{eq:4-1}
\end{equation}

其中 $\alpha$、$\beta$、$\gamma$ 为权重系数，默认值分别为 0.3、0.3、0.4。Importance 的权重略高（0.4），因为其综合了四个独立信号源，在信息多样性和可靠性上均优于单信号的 Recency 和 Uniqueness。以下逐一阐述三个分项的设计动机和计算方法。

\subsection{新近性（Recency）：基于艾宾浩斯遗忘曲线的时间衰减}

新近性刻画了"越近的事件越可能被追问"这一直觉。其数学模型直接借用了艾宾浩斯遗忘曲线的基本形式——指数衰减：

\begin{equation}
\text{Recency}(n) = e^{-\lambda \cdot \Delta t}
\label{eq:4-2}
\end{equation}

\begin{equation}
\lambda = \ln(2) / \textit{half\_life}
\label{eq:4-3}
\end{equation}

\begin{equation}
\Delta t = (\text{当前时刻} - n.\textit{timestamp}).\text{total\_seconds}()
\label{eq:4-4}
\end{equation}

其中半衰期（half\_life）为 Recency 衰减到 0.5 所需的时间间隔，是一个可根据部署场景灵活配置的参数。对于以"天"和"周"为基本回溯单位的长期机器人记忆，建议将半衰期设置为与典型问答时间尺度匹配的值——例如半衰期设为 1 天（86400 秒），此时 1 天前的事件 Recency = 0.50，1 周前的事件 Recency $\approx$ 0.008。Recency 以指数速率平滑衰减：近期事件（同一 episode 内部，相隔数分钟至数小时）之间可获得清晰的区分度（如 5 分钟前 Recency $\approx$ 0.93，1 小时前 Recency $\approx$ 0.97）；中期事件（数天至数周前）的 Recency 逐渐降至较低水平，为后续由 Uniqueness 和 Importance 主导的效用评估提供渐进过渡；远期事件（数月前）的 Recency 已非常微弱但不会对效用函数的整体行为产生支配性影响。

这种指数衰减的设计使 Recency 在效用函数中扮演了"平滑过渡器"的角色：它不是简单地将事件划分为"近期"和"远期"两个离散类别，而是提供了一条连续的时间衰减谱线，使效用函数从以 Recency 为锚点到以 Importance 和 Uniqueness 为锚点的过渡是渐进的。这与人类记忆的双过程理论（dual-process theory）形成有意义的平行：近期记忆依赖海马体介导的情景回忆，对时间上下文高度敏感；远期记忆更多依赖新皮层介导的语义整合，对知识的总体结构和关联性敏感。在 GRAF-Mem 中，Recency 模拟了前者的时间敏感性——近期经历享有最高的综合效用因而最不容易被压缩；Importance 和 Uniqueness 模拟了后者的结构敏感性——远期经历中只有那些在任务网络中处于结构枢纽位置、携带关键用户对话或异常事件的节点才会被优先保留。

\subsection{独特性（Uniqueness）：语义不可替代性}

独特性衡量一个事件节点在其兄弟节点（同一 Goal 下的其他 L2 事件）中的语义不可替代性。其计算基于文本嵌入向量：

\begin{equation}
\text{Uniqueness}(n) = 1 - \text{mean}_{s \in \textit{siblings}(n)} \textit{cos\_sim}(\textit{emb}(n), \textit{emb}(s))
\label{eq:4-5}
\end{equation}

其中 $\textit{emb}(n)$ 为节点 $n$ 的索引内容文本经 all-mpnet-base-v2 模型编码后的 768 维嵌入向量。若 $n$ 没有兄弟节点（$N=1$），Uniqueness 直接取 1.0。

设计直觉为：在兄弟节点中语义越独特的事件，提供的信息越不可由其他节点替代，因此越应被保留。例如，在同一 Goal 的 20 个连续导航帧中（全都是 "Forward()" + 相似的视觉观察），每个帧的 Uniqueness 接近于 0——因为删除其中任何一个，其信息几乎可被周围的兄弟帧完整替代。反之，一个包含用户语音（ASR 文本）的事件或一个动作状态从 success 突变为 failure 的事件，其文本嵌入与周围的重复导航帧显著不同，Uniqueness 值高，受到保护。

值得注意的是，本文采用连续独特性分数而非离散的"重复/不重复"分类——不需要设置阈值判断"两个事件是否足够相似可以合并"，只需要计算"这个事件和周围事件的平均差异度"。这避免了离散合并操作对树结构的侵入式修改（不需要实际合并节点或重写树结构），使遗忘成为一个纯粹的\emph{信息削减}操作而非\emph{结构重写}操作。

\subsection{重要性（Importance）：多信号综合评分}

重要性是效用函数中信息维度最丰富的分项，综合了四个独立的信号源：

\begin{equation}
\text{Importance}(n) = \textit{signal\_asr}(n) + \textit{signal\_failure}(n) + \textit{signal\_goal}(n) + \textit{signal\_centrality}(n)
\label{eq:4-6}
\end{equation}

各信号的具体定义为：

\textbf{（a）用户对话信号（signal\_asr）：}若节点的 ASR 识别文本非空（即该时刻有用户的语音输入），加 0.35。用户语音交互——如任务指令（"make coffee"）、指引反馈（"the bread is in the lower cabinet"）、确认对话（"good job, task complete"）——是回答"当时发生了什么"的最直接线索。任何包含此类语音的节点都受到显著的保留保护。

\textbf{（b）动作异常信号（signal\_failure）：}若节点的动作状态为非成功（不以 "succe" 开头——即动作结果为失败或中断），加 0.20。失败经历在诊断性问答中（"为什么那次任务没有完成""之前遇到过类似的错误吗"）具有极高的信息价值，且失败事件的出现频率远低于成功事件（在 TEACh 数据中通常 $< 5\%$），保护它们的存储代价很小。

\textbf{（c）目标状态变化信号（signal\_goal）：}若节点的目标状态字段非空（即该时刻发生了子目标的完成或切换），加 0.10。目标状态变化标记了任务的起止边界——这些边界事件在按时间线回答"完成了哪些任务"时是关键的锚点。

\textbf{（d）图度中心性信号（signal\_centrality）：}利用事件关联图的度中心性来评估节点的结构重要性。$\textit{signal\_centrality} = 0.35 \cdot \min(\textit{degree}(n) / \textit{max\_degree}, 1.0)$，其中 $\textit{degree}(n)$ 为节点 $n$ 在临时轻量图中的度数（连接的邻居数），$\textit{max\_degree}$ 为当前图中所有节点的最大度数。该信号的设计直觉为：\emph{在事件图中连接越多的节点，作为"结构枢纽"的价值越高}——即使该节点本身的语义独特性不高或时间较久远，但它连接了大量其他事件，在检索中提供关键的横向桥接作用。

由于完整的记忆图（第三章）的构建包含嵌入计算（相似动作边），直接复用完整图来计算中心性会引入额外的开销。因此，遗忘模块构建一个\emph{临时轻量图}——仅包含时间相邻边和共享物体边两种零成本规则边——专门用于度中心性的快速近似计算，而非复用完整记忆图。在 $|h|=50$ 的设置下，临时图的构建时间约 1–3 秒。

Importance 信号的四个分项来源独立、信息互补，且不涉及 LLM 调用，计算成本低。最后一个值得注意的设计决策是 Importance 的钳位——各分项加和后若超过 1.0 则钳位为 1.0，确保单个节点即使满足了所有保护条件（同时有对话+失败+目标变化+高度数），也不会产生溢出效应影响整体效用的归一化。

\section{渐进式遗忘策略}

\subsection{三级遗忘操作}

在计算每个 L2 事件节点的效用值 $U(n)$ 后，系统根据两个阈值 $\theta_1$ 和 $\theta_2$（$\theta_1 > \theta_2$）将节点划分为三个等级，每个等级执行不同程度的压缩：

\textbf{Level 0（完整保留，$U \ge \theta_1$）。}事件的所有层级数据——L0 原始感知（图像、音频）、L1 场景图（物体列表和关系）、L2 摘要文本、L3 目标关联——全部保留。Level 0 节点在存储和检索行为上与未压缩的记忆树完全一致，支持全部问答类型（包括依赖图像帧的 VQA）。

\textbf{Level 1（去细节化，$\theta_2 \le U < \theta_1$）。}删除 L0 层的原始感知数据——将每个关联场景的各帧的原始图像数据和原始音频数据清空。保留全部文本字段：L1 场景图中的物体列表和关系文本、L2 事件摘要、L3 目标摘要。由于图像和音频占据了单个事件存储空间的绝大部分（在 TEACh 数据中，一帧 RGB 图像约 200KB，而该帧的场景图文本约 2KB），Level 1 压缩可节省约 80\% 的单节点存储，但对纯文本问答的影响微乎其微——因为 Agent 在绝大多数问题（除显式要求看图回答的 VQA 外）中仅使用文本摘要和场景图关系。Level 1 是"以最小的功能损失换取最大的存储收益"的甜点区。

\textbf{Level 2（摘要保留，$U < \theta_2$）。}仅保留事件的自然语言摘要文本和时间戳。具体操作包括：（a）将所有场景图缩减至最后一个，并清空其物体列表、关系列表、原始图像和原始音频；（b）将自然语言摘要缓存为摘要覆盖属性，保证检索和展示的一致性；（c）清空音频描述、动作参数摘要等辅助文本字段。Level 2 节点的存储仅剩一串自然语言摘要文本（约 200–500 字符）和一个时间范围——占原始节点存储的约 5\%——支持 overview/时间线/任务结果类问题的回答，但不支持对象邻接、精确位置和局部 relation 类推理（因为这些底层信息已被清空）。

从认知科学的角度来看，这三级操作映射了人类记忆中"情景记忆（episodic memory）→ 语义记忆（semantic memory）→ 要旨痕迹（gist trace）"的渐进式转化谱系。Level 0 对应于生动的情景记忆（保留了完整的多模态感知细节和时空上下文）；Level 1 对应于去情境化的语义记忆（剥离了具体的图像和音频痕迹但保留了结构化的文本知识）；Level 2 对应于模糊的要旨痕迹（仅保留最高层的语义概括）。这种映射使遗忘模块不仅具有工程合理性，也具有认知科学性——它模拟的不是随机的数据删除，而是人类记忆系统在长时程中自然发生的记忆表征转化过程。

三个等级的划分并非按时间窗口（如"删除 30 天前的所有数据"）的硬切分，而是由效用函数连续打分和两个阈值的软划分决定。这意味着：（1）阈值 $\theta_1$ 和 $\theta_2$ 可根据部署场景的存储预算灵活调整（存储压力大时升高阈值以压缩更多节点）；（2）即使是很久以前的事件，只要其重要性高（如失败的诊断性对话），仍可保持在 Level 0 或 Level 1；（3）近期但高度重复的冗余事件也可以被压缩到 Level 1 甚至 Level 2。

\begin{table}[htbp]
\centering
\caption{三级遗忘操作概览}
\label{tab:4-1}
\begin{tabular}{|p{1.8cm}|p{2.5cm}|p{3.5cm}|p{3.5cm}|p{2cm}|p{2.5cm}|}
\hline
等级 & 条件 & 保留内容 & 删除/清空内容 & 典型单节点空间节省 & 支持的问答类型 \\
\hline
Level 0 & $U \ge \theta_1$ & 全部 L0–L4 数据 & 无 & 0\% & 全部（含 VQA） \\
\hline
Level 1 & $\theta_2 \le U < \theta_1$ & 全部文本（场景图、摘要、目标） & L0 图像/音频 & $\sim$80\% & 纯文本问答 \\
\hline
Level 2 & $U < \theta_2$ & 自然语言摘要文本 + 时间戳 & 场景图数据、关系列表、物体列表、所有辅助字段 & $\sim$95\% & overview / 时间线 / 任务结果 \\
\hline
\end{tabular}
\end{table}

默认阈值设置为 $\theta_1=0.5$、$\theta_2=0.2$。在此设置下，一个同时满足"含用户对话（+0.35）+ 动作异常（+0.20）"的节点即使没有图中心性且 Recency 和 Uniqueness 均为 0，其效用值为 $0.4 \times 0.55 = 0.22$——恰好超过 $\theta_2=0.2$，落入 Level 1 而非 Level 2。这是有意而为的：失败对话经历即使时间久远（Recency$\approx$0）且属于常见失败模式（Uniqueness$\approx$0），仍应保留完整的文本信息。

\subsection{遗忘豁免与安全机制}

在上述效用函数和阈值划分的基础上，系统设置了四重保护机制，确保关键证据在任何参数配置下都不被误删：

\textbf{（1）对话豁免。}包含用户语音（ASR 文本非空）的事件节点被标记为遗忘豁免状态，在效用计算和阈值划分后被强制提升至不低于 Level 1（即永远不会进入 Level 2 去细节化）。这是最高优先级的保护规则——用户对话是回答"当时发生了什么"的最直接且唯一的线索，丢失后无法从其他节点恢复。

\textbf{（2）失败豁免。}动作状态为失败或中断的事件节点同样获得遗忘豁免标记，享有与对话豁免同等的保护（不低于 Level 1）。失败经历在事件总体中的占比通常 $< 5\%$，保护它们的存储代价微不足道，但其在诊断性问答中的信息价值远超任何其他节点。

\textbf{（3）Goal 边界保护。}每个 Goal（L3 目标摘要）下的首事件和尾事件被标记为不能进入 Level 2——即至少保持在 Level 1。首事件记录了一个子目标开始时的状态和初始操作（通常是收到指令或开始搜索物体），尾事件记录了子目标完成时的状态和确认（通常是收到用户确认或 placement 完成）。这两个"锚点"在任务 list 类和时间线类问答中提供关键的时间边界标记。

\textbf{（4）安全下限机制。}系统设置最低保留比例参数（默认 0.3），表示至少 30\% 的事件必须保持在 Level 0。在阈值划分完成后，若实际 Level 0 比例低于最低保留比例参数，系统自动下调 $\theta_1$ 和 $\theta_2$（每次步进 0.05），然后重新划分，直至 Level 0 比例满足下限。该机制防止了因阈值设置不当或效用分布异常导致的灾难性遗忘——即使所有事件的效用都极低（如所有事件都是数月前的重复导航帧），系统也会强制保留至少 30\% 的完整事件。

\subsection{遗忘→图构建的先筛后连管线}

遗忘模块与事件关联图模块（第三章）之间的接口设计遵循"先筛后连"的原则：记忆巩固（遗忘）在调用顺序上严格先于事件关联图构建。具体地，在系统初始化函数的执行顺序中，遗忘配置相关的记忆巩固操作在图配置相关的图构建操作之前被弹出和执行。

这一顺序的意义在于：遗忘先精简了树中的低效用节点（大量的重复导航帧被压缩为 Level 1 甚至 Level 2，其底层物体和关系被清空），图在精简后的节点集上构建——节点数不变（遗忘不会删除节点），但 Level 2 节点的物体集合和位置签名字段被清空，使得它们无法参与共享物体边和共位置边的构建。结果：图更小、边更稀疏、检索更聚焦于高价值证据。此外，临时轻量图（用于中心性计算的共享物体边）同样受益于此——精简后的边计算更快、更准确。

"一图两用"是本模块的核心创新点：事件图不仅服务于检索（第三章），其中心性信号也直接进入遗忘的效用函数（Importance 分项中的 signal\_centrality）。通过这种"检索用边做扩展、遗忘用中心性做保护"的协同设计，图结构成为贯穿 GRAF-Mem 多个模块的统一信息基础设施——而非某个模块的附属品。

\section{TEACh 遗忘实验}

\subsection{实验设置}

遗忘实验在 TEACh $|h|=50$ 的 10 段长历史（共 100 QA）上进行，设置两档实验配置：

\textbf{Forget+Graph（保守配置）。}使用默认效用权重（$\alpha=0.3$, $\beta=0.3$, $\gamma=0.4$）和默认阈值（$\theta_1=0.5$, $\theta_2=0.2$），启用图中心性（\texttt{use\_graph\_centrality=True}）。在该配置下，遗忘主要发生在 Level 1（去图像），少量低效用节点进入 Level 2。这是遗忘模块在长期部署中的保守稳定设置——压缩温和，且通过 full 100/100 的答案逐条一致性验证了安全性。

\textbf{Level 2 / Summary-Retention（强压缩配置）。}下调阈值以扩大 Level 2 的比例（$\theta_1 \approx 0.4$, $\theta_2 \approx 0.15$），同时采用修正后的 Level 2 表示——清理缓存的摘要文本（避免同一段摘要文本被序列化两次）、清空残留的音频描述字段和动作参数摘要字段、确保摘要覆盖属性的读取路径与检索索引一致。该配置旨在验证遗忘模块的强压缩潜力——将大量低效用旧记忆真正降级为"仅保留摘要"。

\textbf{评测指标。}主指标为存储压缩的相对比率——树结构缓存大小比、scene 数比、relation 数比（均以未压缩的 Base 设置为 1.0000 进行归一化）。需要注意的是，这里的缓存文件仅包含树结构元数据和文本摘要（不含原始图像/音频文件），因此压缩比反映的是遗忘机制对树结构本身的精简效果；在实际部署中，Level 1 删除 L0 原始感知数据和 Level 2 清空底层关系的操作将产生远高于树结构压缩比的绝对存储节省。准确率方面，使用 QA 答案文本逐条一致性而非 S\textsubscript{c}/S\textsubscript{p} 的自动 Judge 分类——原因在于自动 Judge 的分类结果受答案表述和 LLM 非确定性的影响，当遗忘前后的 raw answer 文本完全一致时，最可靠的结论是"问答行为未发生变化"。

\subsection{存储压缩主结果}

\begin{table}[htbp]
\centering
\caption{TEACh $|h|=50$ 遗忘模块存储压缩与答案一致性}
\label{tab:4-2}
\begin{tabular}{|l|c|c|c|c|c|}
\hline
Setting & QA 范围 & Pickle ratio & Scene ratio & Relation ratio & QA 答案一致性 \\
\hline
Base（无遗忘） & 100/100 & 1.0000 & 1.0000 & 1.0000 & 正式基线 \\
\hline
Forget+Graph & 100/100 & \textbf{0.9474} & 1.0000 & 1.0000 & \textbf{100/100 逐条一致} \\
\hline
Level 2 / Summary-Retention & 100/100 & \textbf{0.6830} & 0.9288 & 0.4917 & \textbf{100/100 逐条一致} \\
\hline
\end{tabular}
\end{table}

从表 4-2 可以得出：

\textbf{（1）保守遗忘的安全性与有限压缩。}Forget+Graph 配置以约 5.26\% 的序列化压缩率，在 full 100/100 QA 上实现了与 Base 完全一致的最终答案（100/100 逐条一致）。场景数和关系数均无任何减少（1.0000）——压缩完全来自 Level 1 操作（删除 L0 图像和音频），说明该配置下所有节点均未进入 Level 2。5.26\% 的压缩率不高，但它证明了核心安全主张：\emph{遗忘可以在不改变任何 QA 输出的前提下压缩存储}。对于真实部署中的保守策略而言，这是一个可以无风险启用的 baseline。

\textbf{（2）强压缩的巨大潜力和答案稳定性。}Level 2 配置将树结构缓存大小压缩至基线的 68.30\%（压缩率 31.70\%），同时将关系数从 577,052 条削减至 283,721 条（削减 50.83\%），场景数从 68,214 降至 63,355（削减 7.12\%）。事件数保持不变（60,836）——遗忘不会删除任何节点，仅对其内容施加不同程度的压缩。即使在这种强度的压缩下，full 100/100 QA 的最终答案仍然与 Base 逐条完全一致。需要强调的是，此处的压缩对象是评测用的树结构元数据（文本摘要、场景图关系等），在实际部署中叠加原始多模态文件后，Level 1 删除 L0 图像/音频和 Level 2 清空底层关系的操作将带来远超 31.70\% 的绝对存储削减——因为原始多模态数据的体量比文本元数据大数个数量级。

\textbf{（3）答案一致性作为更可靠的评测指标。}在 Forget+Graph 的 100/100 答案与 Base 完全一致，但同一批答案的自动 Judge 分类结果却报告了不同的 S\textsubscript{c} 值（45.0\% vs 48.0\%）。这不是遗忘导致了质量下降——相同文本给出不同 judge 分类，只能说明 judge 本身受批次效应和模型非确定性影响。因此本文在遗忘实验中使用 raw answer 逐条一致性作为准确率稳定性的主证据。

\begin{table}[htbp]
\centering
\caption{Level 2 强压缩的详细结构变化}
\label{tab:4-3}
\begin{tabular}{|l|c|c|c|}
\hline
指标 & Base（无遗忘） & Level 2 后 & 变化 \\
\hline
事件数（Events） & 60,836 & 60,836 & 0\%（无事件被删除） \\
\hline
场景数（Scenes） & 68,214 & 63,355 & $-$7.12\% \\
\hline
关系数（Relations） & 577,052 & 283,721 & \textbf{$-$50.83\%} \\
\hline
树结构缓存大小（相对比率） & 1.0000 & 0.6830 & \textbf{$-$31.70\%} \\
\hline
\end{tabular}
\end{table}

场景数的 7.12\% 下降来自于 Level 2 操作中的场景缩减（多个场景合并为最后一个）。关系数的 50.83\% 大幅下降来自两个因素的叠加：Level 2 节点的场景被清空 relations 字段（直接删除），以及 Level 2 节点的 objects 被清空后无法再参与图构建中的共享物体边（间接影响）。树结构缓存 31.70\% 的降幅小于 relations 的降幅，是因为 summary 文本和树结构元数据（时间戳、层级关系、引用指针等）不受遗忘影响，它们构成了缓存的固定基础开销。需要再次强调的是，以上数字仅反映树结构元数据的压缩效果——在实际部署中，Level 1 和 Level 2 操作删除的原始图像和音频文件将贡献远高于树结构压缩的绝对存储节省量。

\subsection{遗忘等级分布}

在 Level 2 强压缩配置下，60,836 个事件的遗忘等级分布为：

\begin{itemize}
\item \textbf{Level 0（完整保留）}：7,071 个事件，占 11.62\%
\item \textbf{Level 1（去细节化）}：16,623 个事件，占 27.33\%
\item \textbf{Level 2（摘要保留）}：37,142 个事件，占 \textbf{61.05\%}
\end{itemize}

超过 61\% 的事件进入 Level 2，说明这轮实验真正触发了遗忘模块的强压缩机制——不是停留在 Level 1 的保守去图像，而是将大部分低效用事件降级为仅保留高层摘要。这一分布与真实长期部署的预期行为一致：在数周至数月的机器人运行中，日常的、重复的、低信息量的操作将构成事件的主体，而关键事件（对话、失败、任务起止）仅占少数。Level 2 机制使得长期存储的增长从 $O(N)$（所有事件的全细节保留）趋近 $O(\log N)$（仅保留高层摘要，其增长速度远低于底层细节）——这是解决长期存储膨胀的关键所在。

\section{ARMARX 遗忘机制验证}

TEACh 的实验证明了遗忘在存储压缩与答案稳定性之间的权衡可以在模拟数据上同时达成。然而，模拟环境中的感知和交互模式与真实机器人存在差异——模拟器中的视觉场景图构建高度结构化，而真实机器人的感知噪声、光照变化和物体杂乱摆放使得事件内容的质量分布更不均匀。ARMARX\_LT\_MEM 数据集（第二章 2.7.2 节）提供了在真实机器人（KIT ARMAR-7）上验证遗忘机制的补充视角。由于该数据集缺少完整的标准化 QA 评测集（仅含 30 对 QA），本节不试图报告完整的 QA 准确率对比，而是通过结构审计和局部细节切片来回答两个更具体的问题：（1）不同遗忘设置在结构层面分别删除了什么？（2）结构层面的差异是否会传递到 relation-sensitive 的问答中？

\subsection{结构审计：遗忘在不同设置下分别删除了什么}

选取 ARMARX\_LT\_MEM 中 117 个高细节密度事件（具有丰富的物体列表和空间关系），分别在五种遗忘设置下统计其结构保留状态：

\begin{itemize}
\item \textbf{base}：无遗忘，全部 117 个事件完整保留
\item \textbf{random\_medium}：随机选择约 50\% 事件施加中等强度遗忘（不基于任何效用信号，纯粹随机基线）
\item \textbf{medium}：普通效用遗忘（$\alpha=0.3$, $\beta=0.3$, $\gamma=0.4$，但 $\gamma$ 中不包含图中心性——\texttt{use\_graph\_centrality=False}）
\item \textbf{medium\_graph}：效用遗忘 + 图中心性（\texttt{use\_graph\_centrality=True}）
\item \textbf{ultra}：超激进遗忘（$\theta_1=0.7$, $\theta_2=0.4$——更高的阈值意味着更多的节点被推向低等级）
\end{itemize}

\begin{table}[htbp]
\centering
\caption{ARMARX 117 个高细节密度事件的结构保真审计}
\label{tab:4-4}
\begin{tabular}{|l|c|c|c|c|c|}
\hline
Setting & L0 & L1 & L2 & Mean scenes & Mean relations \\
\hline
base & 117 & 0 & 0 & 47.31 & 145.44 \\
\hline
random\_medium & 102 & 6 & 9 & 45.01 & 137.20 \\
\hline
medium & 71 & 23 & 23 & 39.74 & 112.63 \\
\hline
\textbf{medium\_graph} & 71 & \textbf{36} & \textbf{10} & \textbf{43.89} & \textbf{130.22} \\
\hline
ultra & 66 & 27 & 24 & 39.58 & 111.82 \\
\hline
\end{tabular}
\end{table}

从表 4-4 中可以获得三个关键观察：

\textbf{（1）medium\_graph 与 medium 的 L0 事件数相同（71），但 L2 事件数显著更少（10 vs 23），且 mean relations 明显更高（130.22 vs 112.63）。}这两个设置对相同数量的事件保留了完整的 Level 0 状态，但 medium\_graph 将更多事件维持在了 Level 1（36 vs 23）而非进一步压缩到 Level 2。在事件图中连接度高的事件（作为多个事件的共享物体或共位置枢纽）因其图中心性获得了额外的重要性加分，从而从 $\theta_2$ 的"引力范围"中被拉回至 Level 1 区间。这验证了图中心性不是简单地"让遗忘更保守"（两个设置的 L0 完全相同），而是精确地"改变遗忘的分布"——在总保留量相同的前提下，优先将结构枢纽事件保留在更高的信息等级。

\textbf{（2）ultra 的 L0 最少（66），但 L2 却不是最多（24 vs medium 的 23）。}超激进遗忘的高阈值将更多事件从 Level 0 推向了两个方向：一部分进入了 Level 1（27 vs medium 的 23），另一部分进入了 Level 2（24 vs medium 的 23）。但对比 medium\_graph 的 10 个 L2 事件，ultra 的 24 个几乎是其 2.4 倍。ultra 的超激进压缩在没有图中心性引导的情况下，会"无差别攻击"——不仅压缩了低价值的重复事件，也压缩了高结构重要性的枢纽事件。

\textbf{（3）random\_medium 在 L0 保留数上最高（102），但它的表现高度不稳定。}随机遗忘在统计意义上保留得更多（因为不区分事件的价值），但实际被遗忘的具体事件完全取决于随机种子——可能保留了大量低价值重复帧而丢失了关键的关系密集事件。这为 4.5.2 的局部细节切片结果埋下了伏笔。

\subsection{局部细节切片：结构差异是否影响问答}

为了验证结构审计中观察到的差异是否会实际传递到问答中，本节构造了 4 个 relation-sensitive 的局部探针（probe）。每个探针仅向模型提供与问题局部相关的最新 scene graph 文本（约 700–1000 字符），而非完整的记忆树。这种设计的目的是在极低 token 成本下隔离观察不同遗忘设置在 relation-sensitive 问题上的纯结构效应——排除检索策略、Agent 导航路径和其他混淆因素。

4 个探针的设计如下：（1）\texttt{objects5\_neighbor\_orangejuice}：询问某物体附近是否出现过 "OrangeJuice"；（2）\texttt{softcake\_counter\_1437}：询问 "softcake" 是否在 "Counter" 附近出现过；（3）\texttt{rusk\_sink\_1456}：询问 "rusk" 是否在 "Sink" 附近出现过；（4）\texttt{armar7\_infront\_counter\_noon}：询问 ARMAR-7 在 "Counter" 前方的某次具体交互。

\begin{table}[htbp]
\centering
\caption{ARMARX 局部细节切片结果（4 个 relation-sensitive probe）}
\label{tab:4-5}
\begin{tabular}{|l|c|c|}
\hline
Setting & 命中数 & 表现描述 \\
\hline
base & \textbf{4/4} & 完整记忆保留了所有局部关系证据 \\
\hline
random\_medium & 3/4 & 不稳定——取决于随机种子是否保留了对应该探针的事件 \\
\hline
medium & 0/4 & 多个 key events 被压成 Level 2，relations 消失 \\
\hline
\textbf{medium\_graph} & \textbf{4/4} & 图中心性保护了 relation-rich 关键事件 \\
\hline
ultra & 0/4 & 过激压缩导致所有局部关系证据丢失 \\
\hline
\end{tabular}
\end{table}

从表 4-5 可以得出，medium\_graph 是唯一一个在保持与 medium 相同 L0 保留数（71）的前提下达到与 base 相同 hit rate（4/4）的设置。这直接验证了 4.5.1 中结构审计的发现：medium\_graph 通过图中心性将 relation-rich 的关键事件维持在 Level 1（保留了 relations 文本），而 medium（无中心性引导）将这些事件推向了 Level 2（清空了 relations），导致 relation-sensitive 问题无法回答。

更值得注意的是，\emph{"忘什么"比"忘多少"更重要}。random\_medium 在 L0 保留数上最高（102/117），但在局部细节切片上只有 3/4 且不稳定——因为它无法区分"哪些事件包含对后续问答关键的 relation 信息"。medium\_graph 在相同甚至更低的 L0 保留数下（71/117），通过精确的选择性遗忘，将有限的保留预算分配给了真正重要的事件。

\section{讨论}

综合 TEACh 的存储压缩实验和 ARMARX 的机制验证实验，本章提出的效用引导渐进式遗忘形成了以下三个核心结论：

\textbf{（1）图中心性引导遗忘实现了从"删多少"到"选择性遗忘"的质变。}表 4-4 和表 4-5 共同构成了最有力的证据链：medium 和 medium\_graph 在总体压缩量上完全相同（L0 均为 71），但 medium\_graph 通过图中心性将 Limited 的 Level 0 保留预算更精准地分配给了 relation-rich 和结构枢纽事件——结果是在局部细节切片上从 0/4 跃升至 4/4。这说明图结构不仅提升了检索质量（第三章），也实质性地提升了遗忘决策的智能化水平。

\textbf{（2）遗忘模块的核心目标：存储压缩与问答稳定性的兼顾——在 TEACh full 100/100 上得到了正式验证。}Forget+Graph 和 Level 2 强压缩两档实验在 full 100/100 QA 上均保持最终答案与 Base 逐条完全一致。这回应了本章引言中提出的 RQ2：在保持问答能力基本稳定的前提下控制存储开销是可行的，且可以通过调节 $\theta_1$ 和 $\theta_2$ 在压缩率和稳定性之间取得不同档位的折中。

\textbf{（3）认知科学的视角：渐进式遗忘模拟了人类记忆表征的自然转化。}61.05\% 的事件进入 Level 2 不是数据的"丢失"，而是记忆表示的"降级"，从多模态细节记忆转变为自然语言摘要记忆。科学文献将人类情景记忆随时间的这种转化描述为"语义化"过程\cite{ref_semantization}：高度情境化的情景记忆痕迹（包含丰富的感知细节和时空上下文）逐渐丧失其知觉细节，转而由去情境化的语义知识（要旨、事实、概括）来支撑回忆。本文提出的 Level 0→Level 1→Level 2 三级遗忘谱系，在工程层面模型了这一认知过程，Level 0 的多模态细节对应于海马体依赖的生动情景回忆，Level 1 的去图像化结构文本对应于部分巩固的混合表征，Level 2 的纯摘要对应于新皮层介导的语义化要旨。Level 2 机制使长期存储的增长从 $O(N)$（所有事件的全细节保留）趋近 $O(\log N)$（仅摘要的递归概括），这是解决 append-only 记忆树在数月甚至数年部署中存储膨胀问题的根本途径。更重要的是，这一方案的科学基础——选择性保留而非均匀遗忘、渐进式降级而非二元删除——使它区别于简单的缓存淘汰策略，为长期机器人记忆管理提供了具有认知科学支撑的系统性框架。
